\section{降低高度却给出错误面积符号的剪切规则}\label{sec:shear}
另一种方案试图通过细分方向来降低高度。它的公式足够初等，可以沿着实际单位线段看清它的吸引力和失败之处。

对单位方向 $u=(c,s)$，当 $cs\ne0$ 时取 $\sigma=\operatorname{sgn}(cs)$，在 $cs=0$ 时保留两种符号。同时应用两个矩阵
\[
 A_\sigma=\begin{pmatrix}1&\sigma\\0&1\end{pmatrix},\qquad
 B_\sigma=\begin{pmatrix}1&0\\\sigma&1\end{pmatrix}.
\]
对其中任意一个矩阵 $M$，令 $m=bu+hn$、$\lambda=|Mu|$、$u'=Mu/\lambda$、$n'=(-u'_2,u'_1)$ 及 $d=(Mm)\cdot u'$。置
\[
 H=h/\lambda,\qquad
 J=[d-(\lambda-1)/2,d+(\lambda-1)/2],\qquad
 \kappa=\operatorname{proj}_{J}(0).
\]
将来源替换为完整单位线段
\begin{equation}\label{eq:shear-unit}
 N'_M=Hn'+(\kappa+[-1/2,1/2])u'.
\end{equation}
撤去所有旧记录；保留两个矩阵以及所有闭接缝分支。将这个固定的来源族操作记为 $F$。

象限的选择保证
\[
 |A_\sigma u|^2=1+2|cs|+s^2\ge1,\qquad
 |B_\sigma u|^2=1+2|cs|+c^2\ge1.
\]
因此，$J$ 非空，式~\eqref{eq:shear-unit} 中的单位线段包含于完整像线段 $MN$ 内。由行列式为一可知，其高度为 $h/\lambda$。对正斜率 $t$，两个分支将 $t$ 分别映到 $t/(1+t)$ 和 $t+1$；它们共同覆盖原来的正象限。负象限按相应的符号规则变化，而保留的闭端点补齐其边界。因此，全方向覆盖得到保留。

逆方向映射从较大的坐标绝对值中减去较小的坐标绝对值。在无理斜率处，减法欧几里得算法使未归一化逆迭代的范数趋于零。沿着传播后的记录追踪，高度因子的乘积逐项消去，因此这些高度趋于零。这确实是已登记高度上的一个有效机制。但它尚未为替换三角形的并集支付面积。

\begin{theorem}[固定双剪切规则的失败]\label{thm:shear-failure}
规则 $F$ 在具有竞争力的全方向来源族上并不保证面积不增。更确切地，存在一列原始有限常数记录输入 $K_j$，满足 $|K_j|\to A_*$，并且
\[
 |F(K_j)|>|K_j|
\]
对每个 $j$ 都成立，其中 $F(K_j)$ 表示按它们指定的来源列表得到的输出。这些正差值不一定有共同的正下界。
\end{theorem}
这里的反例输入将由近极小种子重新构造；结论不声称任意给定的近极小输入本身都会被 $F$ 增大。

\begin{proof}
从面积小于 $\pi/4$ 的有限竞争者出发。它有一条已登记的非零高度单位线段。否则，每个射影方向上的正负径向延伸长度之和都至少为一，极坐标中的平方和不等式就给出面积至少为 $\pi/4$。将一条非零高度记录旋转到水平轴方向，写成
\[
 (b+[-1/2,1/2])e_x+h e_y,\qquad h\ne0.
\]
在这个方向上，$A_+$ 与 $A_-$ 都有 $\lambda=1$。截取区间是单点，因此这两个分支保留高度 $h$，并将中心改为 $b+h$ 和 $b-h$。重复 $k$ 次后，该族包含中心位于下列位置的水平单位线段：
\[
 b+(2j-k)h,\qquad j=0,\ldots,k.
\]
它们的单位底边区间间距为 $2|h|$。实际并集的长度为 $1+k\min(2|h|,1)$：相邻区间或者重叠成一个区间，或者内部不交。从原点对这个并集取锥，所得面积就是其直线长度乘以 $|h|/2$。因此，重复 $k$ 次后的实际体 $K_k$ 满足
\begin{equation}\label{eq:axis-growth}
 |K_k|\ge\frac{|h|}{2}\bigl(1+k\min(2|h|,1)\bigr)\longrightarrow\infty.
\end{equation}
这个估计没有使用带重数的求和。在每个中间高度，水平截面都恰好是同一个区间并集的缩放副本。

令 $a$ 为种子面积，$k$ 为第一个满足 $|K_k|>a$ 的指标。则 $|K_{k-1}|\le a$，而其一步输出的面积更大。这给出一个具有竞争力的紧致来源反例。每个有限次迭代仍然紧致有界，虽然不需要一个对 $k$ 一致的界。

我们仍须恢复原始有限输入，同时保持输出规则不变。固定一次迭代后，在以斜率描述方向的有理坐标图中，传播后的图由有限多个闭的连续半代数片组成。在坐标轴和截取转换处分片，并保留单点片。每个片上的高度要么恒为零，要么其绝对值有正的下界：它等于原始常数高度除以一个连续正值的剪切范数乘积。

对每个方向片作不断加细的闭分割，在每格把参数对 $(b,h)$ 一起采样，再冻结为常数记录。零高度严格保持为零，每个零维片也都原样保留。一致连续性给出完整三角形图的一致逼近。应用 $F$ 后，同样的一致逼近仍然成立，因为它的分支方向映射只依赖方向，而其余公式在每个固定紧分支上一致连续。输入与输出都落在各自紧极限的不断收缩的邻域中，从而得到两个面积上极限估计。

为得到下极限，分开处理非零高度片和零高度片。非零高度片的并集是紧半代数集。其生成三角形内部的并集 $V$ 是半代数的，并在该并集中稠密，因此其边界面积为零。$V$ 中每一点都严格位于某个实际三角形的内部，这个三角形就是它的见证。同一方向上冻结后的三角形最终会包含这个点。Fatou 引理给出面积下极限估计。

在恒为零高度的片上，正负径向延伸长度分别为 $\max(b+1/2,0)$ 和 $\max(1/2-b,0)$。在同样的冻结下，它们一致收敛。因此，在径向边界图之外，这些径向星形体的示性函数几乎处处收敛。应用 $F$ 后，高度仍为零，像方向区间保持固定；相同论证适用于截取后的中心。合并有限多个非退化片和径向片，对冻结输入 $K^{(r)}$ 得到
\[
 |K^{(r)}|\to|K_{k-1}|,
 \qquad |F(K^{(r)})|\to|K_k|.
\]
对足够精细的原始有限输入，严格间隙仍然保留。Hausdorff 逼近提供上极限估计；严格的三角形内部见证以及原样保留的径向片提供下极限估计。

最后，任取一个面积小于 $\pi/4$ 的有限近极小化序列。对每个成员执行同样的旋转和首次越界构造。得到的紧致反例输入，其面积至多为种子面积。选取趋于零的最终有限恢复误差，并使其足够小，以保留该成员的正间隙。按定义，其有限反例输入的面积至少为 $A_*$，并至多为种子面积加上趋于零的恢复误差。这些反例输入的面积趋近 $A_*$，从而证明断言。
\end{proof}

有理轴作为方向标签集合可以忽略，但它保留下来的三角形作为实际材料不能忽略。这就是几乎处处的高度衰减不能控制面积的原因。结论针对的是上面固定的分支与截取规则。改变操作就提出了一个新的几何问题；这个定理并没有把轴上的累积变成所有可能的高度降低操作都会遇到的障碍。

